\chapter*{Kết luận}
\addcontentsline{toc}{chapter}{Kết luận}

\section*{Kết quả đạt được}
\addcontentsline{toc}{section}{Kết quả đạt được}

Khóa luận đã xây dựng được một khung xử lý tuyển dụng số theo hướng cân bằng giữa tính chính xác kỹ thuật và khả năng giải thích nghiệp vụ. Kết quả quan trọng nhất là mô hình hai lớp đánh giá: lớp điều kiện bắt buộc để đảm bảo tính khả thi và lớp tương đồng ngữ nghĩa để xếp hạng mức độ phù hợp. Cách tổ chức này giúp kết quả so khớp có căn cứ rõ ràng, tránh phụ thuộc hoàn toàn vào suy luận khó kiểm soát.

Bên cạnh đó, nghiên cứu đã chuẩn hóa được cách nhìn hệ thống ở mức vận hành: dữ liệu đầu vào được xử lý theo tiến trình trạng thái, tương tác tuyển dụng được theo dõi theo vòng đời, và các sự kiện quan trọng có thể truy vết để phục vụ kiểm tra hoặc cải tiến. Đây là nền tảng cần thiết để chuyển từ bản thử nghiệm sang sản phẩm có khả năng vận hành ổn định.

\section*{Hạn chế}
\addcontentsline{toc}{section}{Hạn chế}

Hạn chế lớn nhất của nghiên cứu nằm ở dữ liệu đánh giá. Khi chưa có tập dữ liệu gán nhãn đủ rộng, việc kết luận chất lượng xếp hạng ở mức tổng quát còn thận trọng. Ngoài ra, dù cơ chế giải thích đã hình thành, cách diễn giải cho người dùng không chuyên vẫn cần được chuẩn hóa sâu hơn để tăng tính nhất quán trong trải nghiệm thực tế.

Một hạn chế khác là độ bao phủ của các tình huống biên trong vận hành. Những trường hợp dữ liệu thiếu, dữ liệu nhiễu cao hoặc biến thể ngôn ngữ đặc thù vẫn cần được quan sát thêm qua các vòng thử nghiệm mở rộng.

\section*{Hướng phát triển}
\addcontentsline{toc}{section}{Hướng phát triển}

Hướng phát triển thứ nhất là xây dựng bộ dữ liệu gán nhãn theo chuyên gia tuyển dụng để đánh giá định lượng có cơ sở hơn. Hướng thứ hai là thực hiện thí nghiệm ablation nhằm đo tác động riêng của từng thành phần trong mô hình, từ đó hiệu chỉnh trọng số theo bằng chứng thay vì kinh nghiệm chủ quan.

Hướng thứ ba là mở rộng cơ chế quan sát vận hành theo thời gian thực để phát hiện sớm sai lệch dữ liệu và suy giảm chất lượng xử lý. Hướng thứ tư là chuẩn hóa ngôn ngữ giải thích cho từng nhóm người dùng, bảo đảm kết quả vừa đúng về kỹ thuật vừa dễ hiểu trong ngữ cảnh nghiệp vụ.

\section*{Kết luận chung}
\addcontentsline{toc}{section}{Kết luận chung}

Tổng thể, khóa luận đã chứng minh tính khả thi của cách tiếp cận tuyển dụng số dựa trên mô hình xử lý có giải thích. Đóng góp của đề tài không nằm ở một thuật toán đơn lẻ, mà ở cách thiết kế hệ thống theo chuỗi xử lý minh bạch, có kiểm soát trạng thái và có khả năng truy vết. Với các hướng phát triển đã nêu, mô hình này có thể tiếp tục mở rộng để đáp ứng yêu cầu thực tế ở quy mô lớn hơn.
